\subsection{二次型群.}

  我们将给 A 上有限维非退化二次型的型所成的集合 \(\mathfrak{W}\) 装备一个交换群结构。我们将用公式在 \(\mathfrak{W}\) 中定义一个加法

\[
\mathrm{T} + \mathrm{T} ^ {\prime} = \theta (\mathrm{T} _ {\tau} \mathrm{T} ^ {\prime})\tag{2}
\]

\[
(\mathrm{T}, \mathrm{T} ^ {\prime} \text {   in   } \mathfrak {W}),
\]

  这个加法是交换的，因为 \(T^{\prime} \tau T\) 等价于 \(T \tau T^{\prime}\)。它是结合的，因为若 T、\(T^{\prime}\) 与 \(T^{\prime\prime}\) 是 W 的元素，我们有

\[
\begin{array}{r l} & (\mathrm{T} + \mathrm{T} ^ {\prime}) + \mathrm{T} ^ {\prime \prime} \sim (\mathrm{T} + \mathrm{T} ^ {\prime}) \tau \mathrm{T} ^ {\prime \prime} \sim (\mathrm{T} \tau \mathrm{T} ^ {\prime}) \tau \mathrm{T} ^ {\prime \prime} \\ & \qquad \sim \mathrm{T} \tau (\mathrm{T} ^ {\prime} \tau \mathrm{T} ^ {\prime \prime}) \sim \mathrm{T} \tau (\mathrm{T} ^ {\prime} + \mathrm{T} ^ {\prime \prime}) \sim \mathrm{T} + (\mathrm{T} ^ {\prime} + \mathrm{T} ^ {\prime \prime}), \end{array}
\]

从而 \((\mathrm{T} + \mathrm{T}' ) + \mathrm{T}'' = \mathrm{T} + (\mathrm{T}' + \mathrm{T}'' )\)，因为 W 中具有相同型的两个元素相等。此外，刚才定义的加法有一个单位元素：事实上，显然所有中性型都有同一个型 \(T_{0}\)，即维数为零的零型的型；立即可见 \(T_{0}\) 就是所求的中性元素。最后，对每个 \(T \in W\) 都存在与 T 相反的元素，这一点立即由下列命题得出：

命题 3. — 设 Q 是 A 上向量空间 V 上一个非退化、有限维的二次型。用 −Q 记 V 上由 \((-Q)(x) = -Q(x) (x \in V)\) 定义的二次型。则型 \(Q \tau (-Q)\) 是中性的。

  事实上，\(Q \tau (-Q)\) 在对角线 \(D\)（\(V \times V\) 的对角线）上的限制是零。因此这个型的指标 \(\geqslant\frac{1}{2}\dim(V\times V)\)（ \(\S\) 4 第 2 小节定义 2），从而等于 \(\frac{1}{2}\dim(V\times V)\)（同出处公式 (4)）。由此得出 \(Q\tau(-Q)\) 是中性的（同出处）。

  这使我们能够陈述下列定义：

定义 1. --- A 上有限维非退化二次型的型所成的集合，带上由 (2) 定义的加法，称为 A 的二次型群，或 A 的维特群。

注. --- 1) 每个型为零（即用上面的记号使 \(\theta(Q) = T_0\)）的非退化、有限维二次型 Q 都是中性型。事实上，存在中性型 N、N′，使得 \(Q \tau N\) 等价于 N′。这表明 Q 是偶维的，从而存在一个与 Q 同维数的中性型 \(N_1\)。由于 Q 与 \(N_1\) 有相同的型，由命题 1 它们等价，故 Q 是中性的。

\begin{enumerate}
\def\labelenumi{\arabic{enumi})}
\setcounter{enumi}{1}
\tightlist
\item
  对 A 上每个有限维二次型 Q，用 \(\delta(Q)\) 记 Q 的维数模 2 的类。我们有
\end{enumerate}

\[
\delta (Q \tau Q ^ {\prime}) = \delta (Q) + \delta (Q ^ {\prime}).
\]

  由于每个中性型 N 都是偶维的，我们有 \(\delta(N) = 0\)；因此关系 \(Q \sim Q'\) 蕴涵 \(\delta(Q) = \delta(Q')\)。于是 \(\delta\) 在 A 上二次型的型所成的群 \(\mathfrak{W}\) 上的限制是 \(\mathfrak{W}\) 到群 \(\mathbf{Z}/(2)\) 中的一个同态。当 A 的特征 \(\neq 2\) 时这个同态是满射的，但当 A 的特征为 2 时则不然，因为这时奇维的二次型是退化的，由于其相伴双线性型是交错的（参见 § 5）。

\begin{enumerate}
\def\labelenumi{\arabic{enumi})}
\setcounter{enumi}{2}
\tightlist
\item
  设 \(a\) 是 A 的一个 \(\neq 0\) 的元素。若 N 是中性型，则 \(aN\) 也是。由此得出关系 \(Q \sim Q'\) 蕴涵 \(aQ \sim aQ'\)。对群 \(\mathfrak{W}\) 的每个元素 T，我们置
\end{enumerate}

\[
a. \mathrm{T} = \theta (a \mathrm{T}).\tag{3}
\]

  这样我们就在 A 的非零元素所成的群 A* 与群 W 之间得到一个外合成法则。下列公式立即由定义得出：

\[
a. (\mathrm{T} + \mathrm{T} ^ {\prime}) = a. \mathrm{T} + a. \mathrm{T} ^ {\prime}, \quad a b. \mathrm{T} = a. (b. \mathrm{T})\tag{4}
\]

\[
(a, b \text {   in   } A ^ {*}, T, T ^ {\prime} \text {   in   } \mathfrak {W}).
\]

  另一方面，若 \(a\)、\(b\) 与 \(a + b\) 都属于 \(\mathbf{A}^*\)，则一般没有 \((a + b)\) . \(\mathrm{T} = a\) . \(\mathrm{T} + b\) . \(\mathrm{T}\)（ \(\mathrm{T} \in \mathfrak{W}\)）。

命题 4. --- 设 Q 是 A 上有限维向量空间 E 上一个非退化二次型。假设 A 的特征 \(\neq 2\)，且设 \((x_{1},\ldots,x_{n})\) 是 V 的一个正交基。用 \(T_{1}\) 记二次型 \(Q_{1}\)（定义在向量空间 A 上且使 \(Q_{1}(1)=1\)）的型。则 Q 的型是 \(\sum_{i=1}^{n}\mathrm{Q}(x_{i}).T_{1}\)。

  事实上，型 Q 等价于

\[
(\mathrm{Q} (x _ {1}) \mathrm{Q} _ {1}) \tau \dots \tau (\mathrm{Q} (x _ {n}) \mathrm{Q} _ {1})
\]

推论. --- 在命题 3 的假设与记号下，元素 \(a\) . \(\mathrm{T}_{1}\)（\(a \in \mathrm{A}^{*}\)）构成 A 上二次型群的一组生成元。

  这样，确定 A 上二次型群的结构就归结为求形如 \(a.T_{1}\) 的元素之间存在的 Z-线性关系。若 \(b \in A^{*}\)，命题 4 中定义的型 \(Q_{1}\) 显然等价于 \(b^{2}Q_{1}\)；因此有 \(a.T_{1} = ab^{2}.T_{1}\)，这表明 \(a.T_{1}\) 只依赖于 a 模子群 \((A^{*})^{2}\)（由 \(A^{*}\) 中元素的平方组成）的类。此外，由命题 3 得出 \((-a).T_{1} = -a.T_{1}\)。然而，除了由刚才指出的关系推出的那些之外，诸 \(a.T_{1}\) 之间一般还存在别的 Z-线性关系。

命题 5. --- 假设 A 是一个极大有序域。设 Q 是 A 上一个非退化、有限维的二次型，(s, t) 是它的符号差（§ 7 第 2 小节定义 2）。则 Q 的型是 (s - t)·T₁，且 A 上二次型的群 W 是由 T₁ 生成的无限循环群。

  事实上，由于 \(\mathbf{A}^{*}/(\mathbf{A}^{*})^{2}\) 的阶为 2，且 \((-1)\) . \(\mathrm{T}_{1} = -\mathrm{T}_{1}\)，故 \(\mathfrak{W}\) 由 \(\mathrm{T}_{1}\) 生成，从而是单演的。对一切 \(n > 0\)，\(n\) . \(\mathrm{T}_{1}\) 是维数 \(n\) 的正非退化二次型的型；由于这些型不是中性的，我们有 \(n\) . \(\mathrm{T}_{1} \neq 0\)，这表明 \(\mathfrak{W}\) 是无限的。最后，用命题 4 的记号，一个符号差为 \((s, t)\) 的型同构于 s 个 \(Q_{1}\) 与 t 个 \(-Q_{1}\) 的直和（§ 7 第 2 小节定理 1）；由此得出它的型是 \((s - t)\) . \(T_{1}\)。

